% !Mode:: "TeX:UTF-8"
% !Mode:: "TeX:UTF-8"
\newcommand{\thepage}{} % to silence a hyperref warning about pdfpagelabels
\RequirePackage[]{silence}
\WarningsOff[hyperref]
\PassOptionsToPackage{log-declarations=false}{xparse}
%\RequirePackage[l2tabu, orthodox]{nag}
\RequirePackage{luatex85} % to make \pdf...= commands work with LuaTeX
\documentclass[aspectratio=169,t]{beamer}
\usepackage{ifthen}
\usepackage[T1]{fontenc}
\usepackage{iftex}
\ifLuaTeX
 \usepackage{fontspec}
 \setmainfont[Ligatures=TeX]{Source Sans Pro}
\else
 \usepackage[utf8]{inputenc}
\fi
\usepackage{german}
\usepackage{newunicodechar}
\newunicodechar{€}{\texteuro}
\newunicodechar{ß}{\ss}
\renewcommand{\SS}{\iffontchar\font"1E9E \symbol{"1E9E}\else SS\fi}
\usepackage{ellipsis}
\usepackage{microtype}
\usepackage{listings}
\usepackage{sourcecodepro}
\lstset{language=Python, basicstyle=\ttfamily}
\usepackage{amsfonts}
\usepackage{amssymb}
\usepackage{amsthm}
\usepackage{fancyhdr}
\usepackage{graphicx}

\newtheorem{proposition}[theorem]{Proposition}

\DeclareMathOperator*{\argmax}{argmax}
\DeclareMathOperator{\supp}{supp}
\newcommand{\E}{\ensuremath{\mathbb{E}}}
\newcommand{\V}{\ensuremath{\mathbb{V}}}

\setbeamercolor{emph}{fg=purple}
\renewcommand<>{\emph}[1]{%
  {\usebeamercolor[fg]{emph}\only#2{\itshape}#1}%
}



\hypersetup{draft=false,colorlinks=false,pdfpagemode=UseOutlines,plainpages=false,
            breaklinks=true,bookmarksnumbered=true,pdfborder=0 0 0,
            pdftitle={Data Science SoSe 2026 HHU},
            pdfauthor={Konrad Voelkel},
            pdfsubject={data science}}

%\documentclass[aspectratio=169,t]{beamer}
\usetheme{hhubeamer} % put the hhu-beamer-theme in /usr/share/texmf/tex/latex/beamer/themes/

% -------------------- Code listings config --------------------
% Define custom colors for syntax highlighting
\definecolor{codegreen}{rgb}{0,0.6,0}
\definecolor{codegray}{rgb}{0.5,0.5,0.5}
\definecolor{codepurple}{rgb}{0.58,0,0.82}
\definecolor{backcolour}{rgb}{0.96,0.96,0.94} % Very light gray/beige
% Configure the listings package
\lstset{
    language=Python,
    backgroundcolor=\color{backcolour},   
    commentstyle=\color{codegreen}\itshape,
    keywordstyle=\color{blue}\bfseries,
    numberstyle=\tiny\color{codegray},
    stringstyle=\color{codepurple},
    basicstyle=\ttfamily\footnotesize, % Use a monospaced font
    breakatwhitespace=false,         
    breaklines=true,      
             % Enable automatic line breaking
    captionpos=b,                      % Caption at the bottom
    keepspaces=true,                 
    numbers=left,                      % Line numbers on the left
    numbersep=10pt,                    % Distance of line numbers from the code
    showspaces=false,                
    showstringspaces=false,            % Don't show weird underscores in strings
    showtabs=false,                  
    tabsize=4,
    frame=single,                      % Add a subtle border around the code
    rulecolor=\color{black!20},        % Border color
  }

\event[]{}
\author[Sommersemester 2026]{Dr. Konrad Völkel, Sommersemester 2026}

%\title{Data Science, VL 1}
%\date{16.04.2026}


%%% Local Variables:
%%% mode: latex
%%% TeX-master: t
%%% End:


\title{Data Science, VL 7} % war in 2025 noch VL 6; Einschub zu Python kam dazwischen...
\date{07.05.2026}

\begin{document}

\begin{frame}
  \titlepage
\end{frame}


\section{Recap}

\begin{frame}
  \frametitle{Recap: Diskrete Zufallsvariablen}
  \begin{itemize}[<+->]
  \item Eine Abbildung $X \colon \Omega \to \Omega'$ mit endlichem oder abzählbar unendlichem Bild $X(\Omega)$ heißt \alert{diskrete Zufallsvariable}, und es gibt eine pmf $p \colon \Omega' \to [0,1]$.\smallskip
  \item Wenn $X \colon \Omega \to \mathbb{R}$, so heißt $X$ \alert{reelle Zufallsvariable} und im diskreten Fall ist $\mathbb{E}X = \sum_{x \in \mathbb{R}} x\ p(x)$ der Erwartungswert und $\mathbb{V}X = \mathbb{E}(X^2) - (\mathbb{E}X)^2$ die Varianz.\smallskip
  \item Im diskreten Fall existieren Erwartungswert und/oder Varianz nicht immer. Wenn $X(\Omega)$ endlich ist, aber schon.\smallskip
  \item Die Entropie von $X$ ist $H(X) = -\sum_{x \in \mathbb{R}} p(x)\log p(x)$, wobei man für die Funktion $x \mapsto x\log x$, in die hier $p(x)$ eingesetzt wird, stets $0 \mapsto 0$ definiert (den Grenzwert für $x\mapsto 0$), sodass $p(x)\log p(x) = 0$, wenn $p(x)=0$ ist.
  \end{itemize}
\end{frame}

\begin{frame}
  \frametitle{Träger einer Zufallsvariable}
  \vspace{-0.5cm}
  \begin{definition}
    Der \alert{Träger} (support) einer Verteilung $P$ auf $\Omega$ ist
    \[\supp(P) := \bigcap \overline{\{ A \subseteq \Omega : P(A) = 1 \}}\]
    (die kleinste [abgeschlossene] Menge, für die $P(\supp P) = 1$ gilt). Für diskrete Verteilungen gilt $\supp(P) = \{x\in\Omega: p(x) > 0\}$.
  \end{definition}
  \medskip
  \uncover<2>{Eine Teilmenge $M \subset \mathbb{R}^n$ heißt \alert{abgeschlossen}, wenn zu jeder konvergenten Folge $(a_n)_{n\in\mathbb{N}}$ in $M$ der Grenzwert $\lim_{n\to\infty} a_n \in \mathbb{R}^n$ auch in $M$ liegt.
  
  Beispiel: abgeschlossene Intervalle $[x, y]$ für $x, y\in\mathbb{R}$ und $\mathbb{R}$ selbst.}
\uncover<3->{
  \vspace{-1.3cm}
  
  Beispiel:\\
    Der Träger der diskreten Gleichverteilung der Zahlen $\{1,\dots,n\} =: \Omega$
    ist $\Omega$ selbst.
  }
  
\smallskip

  \uncover<4->{Beispiel 2:\\
    Die diskrete Gleichverteilung der Zahlen $\{1,\dots,n\}$, aufgefasst als Verteilung auf $\mathbb{R}$ hat den Träger $\{1,\dots,n\}$.}
\end{frame}

\section{Stetige ZV}

\begin{frame}
  \frametitle{Wahrscheinlichkeitsdichte}
  \begin{definition}
    Für ein Wahrscheinlichkeitsmaß $P$ auf $\mathbb{R}$
    heißt eine Funktion $p \colon \mathbb{R} \to \mathbb{R}$
    eine \alert{Wahrscheinlichkeitsdichte} (\alert{pdf, probability density function}),
    wenn für jede messbare Menge $A \subseteq \mathbb{R}$ gilt:
    \[P(A) = \int_A p(x) dx.\]
  \end{definition}
\medskip
  \uncover<2->{Beispiel: Wenn $A=[a,b] \subset \mathbb{R}$ ist, dann ist $P([a,b]) = \int_a^b p(x) dx$.}
\end{frame}

\begin{frame}
  \frametitle{Pdf der stetigen Gleichverteilung}
  \vspace{-1cm}
  \begin{block}{Beispiel}
    Auf dem Intervall $[0,1]$ ist ein Wahrscheinlichkeitsmaß $P$ definiert durch
    die Funktion $p \colon \mathbb{R} \to \mathbb{R}$ mit
    \[p(x) := \begin{cases} 1 & 0 \leq x \leq 1 \\ 0 & \text{sonst.} \end{cases}\]
    D.h. für $A \subseteq [0,1]$ eine messbare Menge ist
    $P(A) := \int_A p(x) dx$
    ein Wahrscheinlichkeitsmaß mit pdf $p$.\pause

    Der Träger ist $\supp P = [0,1] \subset \mathbb{R}$.
  \end{block}

\medskip \pause 
  
Beispiel:
\begin{itemize}
\item Für $A = [a,b] \subseteq [0,1]$ ist $P(A) = \int_a^b 1 dx = [x]_a^b = b-a$.\smallskip\pause
\item Auch $[a,b)$ und $(a,b]$ und $(a,b)$ haben Wahrscheinlichkeit $b-a$.
\end{itemize}
\end{frame}

\begin{frame}
  \frametitle{Nullmengen}
  Ein Integral über $p$ hängt nicht von einzelnen Werten von $p$ ab:\pause

\begin{itemize} 
\item Sei $\tilde{p}$ die Funktion mit $\tilde{p}(x)=p(x)$ für alle $x \neq 0$ und $\tilde{p}(0) = \alpha$ für ein beliebiges festes $\alpha \in \mathbb{R}$. Dann ist $\int \tilde{p}(x) dx = \int p(x) dx$.\smallskip\pause
\item Für die stetige Gleichverteilung auf $[0,1]$ gilt $P(\{x\}) = 0$ für jedes $x \in [0,1]$ und $P(A) = 0$ für jede endliche Menge $A$, obwohl $p(x) \neq 0$.\smallskip\pause
\end{itemize}

\bigskip

\begin{block}{Knobelaufgabe}
Kann es eine Verteilung $P$ auf $[0,1]$ geben mit $P(\{x\}) \neq 0$ für alle $x \in [0,1]$?
\end{block}
\end{frame}

\begin{frame}
  \frametitle{Unendliche Nullmengen}
  Die Menge $\mathbb{Q}$ hat das Lebesgue-Maß $0$, insb. gilt für die stetige Gleichverteilung $P$
  
  \[P(\mathbb{Q} \cap [0,1]) = 0.\]

\bigskip\pause

\begin{itemize}
\item Intuitiv gilt dies, da $\mathbb{Q}$ abzählbar ist, während $[0, 1] \subset \mathbb{R}$ überabzählbar ist.\smallskip\pause
\item Man spricht von einer \alert{Nullmenge}.\smallskip\pause
\item Ein Integral über eine Funktion $f$ ändert sich nicht, wenn man $f$ auf einer Nullmenge verändert.
\end{itemize}
\end{frame}

% \begin{frame}
%   \frametitle{Nicht-messbare Mengen}
%   \begin{theorem}[Satz von Vitali]
%     Für je $x,y \in \mathbb{R}$ schreiben wir $x \sim y$ wenn $x-y \in \mathbb{Q}$ liegt.
%     Diese Relation ist reflexiv, symmetrisch und transitiv, also eine Äquivalenzrelation.
%     Die Äquivalenzklassen $[x] = \{y \in \mathbb{R} : y \sim x\}$ bilden eine disjunkte Zerlegung von $\mathbb{R}$.
%     Sei $V \subset \mathbb{R}$ ein Repräsentantensystem, d.h. eine Menge, die zu jeder dieser Klassen genau ein Element enthält.
%     Dann ist $V$ (die \alert{Vitali-Menge}) nicht Lebesgue-messbar\uncover<2->{, d.h. es lässt sich kein Wert $\lambda(V)$ für die Größe von $V$ festlegen, wenn man für Intervalle $[a,b]$ die Größe $\lambda([a,b]) = b-a \in \mathbb{R}$ festlegt (das nennt man das Lebesgue-Maß $\lambda$).}
%   \end{theorem}
% \end{frame}

% \begin{frame}
%   \frametitle{Nicht-messbare Mengen}
%   \begin{proof}[Beweis von Vitali]
%     Sei $\{q_1,\ q_2,\ \dots\} = \mathbb{Q} \cap [-1,1]$ eine Aufzählung und
%     $V_k := V+q_k$ für $k \in \mathbb{N}$.
    
%     \[[0,1] \subseteq \bigcup_k V_k \subseteq [-1,2]\]
%     somit $1 = \lambda([0,1]) \leq \lambda \left( \bigcup_k V_k \right) \leq \lambda([-1,2]) = 3$.

%     Da die $V_k$ paarweise disjunkt sind,
%     ist $\lambda \left( \bigcup_k V_k \right) = \sum_{k=1}^\infty \lambda \left(  V_k \right)$,\\
%     und $\lambda(V_k) = \lambda(V)$ (denn $\lambda$ ist translationsinvariant).

%     Also ist $1 \leq \sum_{k=1}^\infty  \lambda(V)$, d.h. $\lambda(V) > 0$.\\
%     Außerdem ist $ \sum_{k=1}^\infty  \lambda(V) \leq 3$, d.h. $\lambda(V) = 0$.
%     Widerspruch!\\
%     Also ist $\lambda(V)$ nicht definierbar.
%   \end{proof}
% \end{frame}
%
%\begin{frame}
%  \frametitle{Messbare Mengen}
%  Welche Mengen sind (Lebesgue-)messbar (bzw. Borel-messbar)?
%
%  Mindestens alle (offenen und halboffenen und geschlossenen) Intervalle und alle endlichen Mengen,
%  sowie daraus alle endlichen und abzählbar unendlichen Vereinigungen und Schnitte und Komplemente.
%  
%  \uncover<2->{
%    Wenn $f \colon \mathbb{R} \to \mathbb{R}$ stetig ist, dann sind Urbilder von messbaren Mengen wieder messbar.}
%  
%  \uncover<3->{Formal:\\ Die Borel-$\sigma$-Algebra ist die kleinste, die alle offenen Mengen enthält (daher stetig$\Rightarrow$messbar), die Lebesgue-$\sigma$-Algebra ist die Vervollständigung davon (um alle Teilmengen von Nullmengen).}
%\end{frame}

\begin{frame}
  \frametitle{Integrierbare Funktionen}
  Nicht jede Funktion ist (Lebesgue-)integrierbar. Welche sind es?
  
  \medskip

  \uncover<2->{
  Insbesondere alle Riemann-integrierbaren Funktionen $f \colon \Omega \subseteq \mathbb{R} \to \mathbb{R}$,
  also insbesondere alle stückweise stetigen Funktionen.}

\smallskip

  \begin{definition}<3->
    Jede integrierbare Funktion $f \colon \Omega \to [0, \infty)$ mit $\int_\Omega f(x) dx = 1$ ist \alert{pdf} des Wahrscheinlichkeitsmaßes $P$ auf $\Omega$ mit $P(A) := \int_A f(x)dx$.
  \end{definition}
\end{frame}

\begin{frame}
  \frametitle{Dichte für diskrete Verteilungen}
  Sei $\Omega = \{1,\dots,n\}$ eine endliche Menge und $X \colon \Omega \to \mathbb{R}$ eine Zufallsvariable.
  \begin{itemize}
  \item Dann gibt es eine pmf $p \colon \mathbb{R} \to [0,1]$ \\\pause
  (die außerhalb von $\supp X = X(\Omega) = \{x_1,\dots,x_n\}$ gleich $0$ ist).\pause
  \end{itemize}
  \medskip
  
  Gibt es eine pdf $f \colon \mathbb{R} \to [0,\infty)$?\pause
  \begin{itemize}
  \item Leider ist $\supp X$ eine endliche Menge, also eine Nullmenge, somit ist $\int_{\supp X} f(x) dx = 0$ für jede Funktion $f$.\smallskip\pause
  \item Also kann nicht zugleich $\int_{\supp X} f(x) dx = 1$ sein, d.h. $X$ hat keine pdf.
  \end{itemize}
\end{frame}

\begin{frame}
  \frametitle{Stetige Zufallsvariablen}
  \begin{definition}
    Eine reelle Zufallsvariable $X \colon \Omega \to \mathbb{R}$ heißt \alert{stetige Zufallsvariable},
    wenn $X$ keine diskrete Zufallsvariable ist (also $X(\Omega)$ überabzählbar unendlich ist) und $X$ eine \alert{Dichtefunktion} hat.
  \end{definition}
\end{frame}


% \begin{frame}
%   \frametitle{Allgemeiner wird's schöner}
%   Man kann nicht nur für das Lebesgue-Maß $\lambda$ den Begriff der pdf definieren,
%   i.A. spricht man von Verteilungen $P$, die \alert{absolutstetig} bezüglich eines Maßes $\mu$ sind.

%   \uncover<2->{Auch diskrete Verteilungen sind absolutstetig bezüglich eines `Abzählmaßes' $\mu$ auf endlichen Mengen
%   und in diesem Sinne gibt es einen einheitlichen Formalismus für pmf und pdf.}

%   \uncover<3->{So abstrakt und allgemein brauchen wir es nicht!}
% \end{frame}

\begin{frame}
  \frametitle{Beispiel einer stetigen ZV}
  \vspace{-0.5cm}
  In der Übung haben wir eine Verteilung auf $[5,8) \subset \mathbb{R}$ simuliert.
  Diese Verteilung lässt sich als Zufallsvariable modellieren:
  \[X \colon \{5,6,7\} \times [0,1) \to \mathbb{R},\ (y,z) \mapsto y+z\]
  wobei wir auf $\{5,6,7\}$ die Verteilung mit Wahrscheinlichkeitsvektor $(\frac{4}{7},\frac{2}{7},\frac{1}{7})$ und auf $[0,1)$ die stetige Gleichverteilung wählen.
  
  \smallskip

  \uncover<2->{
    Diese Verteilung hat eine pdf $p \colon [5,8) \to [0,\infty)$ mit
    \[p(x) = \begin{cases} \frac{4}{7} & 5 \leq x < 6 \\ \frac{2}{7} & 6 \leq x < 7 \\ \frac{1}{7} & 7 \leq x < 8 \end{cases} \]
    $P([5,6]) = \int_5^6 p(x)dx = \frac{4}{7} [x]_{x=5}^{x=6} = \frac{4}{7}$
    und $P([6,6.5]) = \int_6^{6.5} p(x)dx = \frac{2}{7} \cdot \frac{1}{2} = \frac{1}{7}$.
    }
\end{frame}

\begin{frame}
  \frametitle{Momente, allgemein}
  \vspace{-0.7cm}
  Man kann für $X \colon \Omega \to \mathbb{R}$ allgemein $\mathbb{E}X$ hinschreiben, was allerdings nicht immer existiert (d.h. eine reelle Zahl ist, also z.B. kann $\mathbb{E}X=\infty$ sein).
  
  \smallskip
  
  Wenn $X$ eine pdf $p \colon \mathbb{R} \to [0,\infty)$ hat, so ist
  \[\mathbb{E}X = \int_{-\infty}^\infty x p(x) dx\]
  das 1-te Moment und weiter
  \[\mathbb{E}\left(X^n\right) = \int_{-\infty}^\infty x^n p(x) dx,\]
  das $n$-te Moment von $X$.
  
  \medskip\pause
  
  Analog ist
  \[\E((X - \E X)^n)\;,\]
  das $n$-te zentrale Moment von $X$.
\end{frame}

\begin{frame}
  \frametitle{Momente der Gleichverteilung}
  Die stetige Gleichverteilung $X$ auf $[0,1]$ (mit pdf die Indikatorfunktion $p = \chi_{[0,1]}$) hat Erwartungswert
  \[\mathbb{E}X = \int_{-\infty}^\infty x p(x) dx\uncover<2->{ = \int_0^1 x dx }\uncover<3->{ = {\left[ \frac{1}{2} x^2 \right]}_0^1} \uncover<4->{ = \frac{1}{2}}\]
  \uncover<5->{und Varianz}
  \[\uncover<5->{\mathbb{V}X = \mathbb{E}(X^2) - {\left(\frac{1}{2}\right)}^2} \uncover<6->{ = \int_0^1 x^2 dx - \frac{1}{4}} \uncover<7->{ = {\left[ \frac{1}{3} x^3 \right]}_0^1 - \frac{1}{4}} \uncover<8->{ = \frac{1}{3} - \frac{1}{4} = \frac{1}{12}.}\]
\end{frame}

\begin{frame}
  \frametitle{Zum Üben: Gleichverteilung auf Intervall}
  Seien $a < b \in \mathbb{R}$ und $p \colon \mathbb{R} \to [0,\infty)$ gegeben durch $p := \frac{\chi_{[a,b]}}{b-a}$, also für $a \leq x \leq b$ ist $p(x) = \frac{1}{b-a}$ und $0$ sonst. Sei $X$ eine stetige ZV mit pdf $p$.

  \[\mathbb{E}X = \uncover<2->{ \int_a^b xp(x) dx = \frac{1}{b-a} {\left[ \frac{1}{2} x^2 \right]}_a^b = \frac{b^2-a^2}{2(b-a)} = \frac{a+b}{2}.}\]
  \[\mathbb{V}X = \uncover<3->{\int_a^b x^2 p(x) dx - \frac{(a+b)^2}{2^2} = \frac{1}{b-a} {\left[ \frac{1}{3} x^3 \right]}_a^b - \frac{(a+b)^2}{4} }\uncover<3>{ = }\]
  \[\uncover<4->{ = \frac{b^3-a^3}{3(b-a)} - \frac{(a+b)^2}{4} = \frac{4b^2 + 4ab + 4a^2 - 3a^2 - 6ab -3b^2}{12} = \frac{1}{3}{\left(\frac{b-a}{2}\right)}^2}\]
\end{frame}

\begin{frame}
  \frametitle{Transformationsformel}
  \begin{lemma}
    Sei $X \colon \Omega \to \mathbb{R}$ eine stetige ZV mit pdf $p \colon \mathbb{R} \to [0,\infty)$ und $f \colon \mathbb{R} \to \mathbb{R}$ eine messbare Abbildung.
    Dann ist
    \[\mathbb{E}(f(X)) = \int_{-\infty}^{\infty} f(x) p(x) dx.\]
  \end{lemma}
\end{frame}

\begin{frame}
  \frametitle{Kumulativ}
  \begin{definition}
  Sei $P$ ein Wahrscheinlichkeitsmaß auf $\mathbb{R}$.
  Die Funktion $F \colon \mathbb{R} \to \mathbb{R}$,
  \[c \mapsto P((-\infty, c))\]
  heißt \alert{kumulative Verteilungsfunktion}, \alert{cdf = cumulative distribution function} von $P$.
\end{definition}
\uncover<2->{
Wenn $P$ eine Dichte $f$ hat, so ist
\[F(c) = P\left((-\infty ,c)\right) = \int_{-\infty}^c f(x) dx\]}
\uncover<3->{
also $F$ eine Stammfunktion von $f$ und $f$ die Ableitung von $F$.}
\end{frame}

\begin{frame}
  \frametitle{Kumulativ}
  \begin{definition}
  Sei $X$ eine reelle ZV.
  Die Funktion $F \colon \mathbb{R} \to \mathbb{R}$,
  \[c \mapsto P(X \leq c)\]
  heißt \alert{kumulative Verteilungsfunktion}, \alert{cdf = cumulative distribution function} von $X$.
\end{definition}
Wenn $X$ eine Dichte $f$ hat, so ist
\[F(c) = P(X \leq c) = \int_{-\infty}^c f(x) dx\]
also $F$ eine Stammfunktion von $f$ und $f$ die Ableitung von $F$.
\end{frame}

\begin{frame}
  \frametitle{Weiter geht's ...}
  ... mit den Notebooks ``Normalverteilung'' und ``wichtigste stetige Verteilungen und Scipy.Stats'' im Skript.
\end{frame}

\begin{frame}[separator]
  \frametitle{Vielen Dank für Ihre Aufmerksamkeit!}
\end{frame}

\end{document}

%%% Local Variables:
%%% mode: latex
%%% TeX-master: t
%%% End:
